Los algoritmos de busqueda son una herramienta fundamental de la Inteligencia Artificial para
resolver problemas en los que un agente debe encontrar una secuencia de acciones que lo lleve
desde un estado inicial hasta un estado objetivo. Cuando el agente no dispone de informacion
adicional sobre la ubicacion de la meta, se habla de \textbf{busqueda no informada}; cuando se
incorpora una funcion heuristica $h(n)$ que estima el costo restante hacia el objetivo, se habla
de \textbf{busqueda informada}.

En este laboratorio se utiliza el entorno de Pac-Man, provisto por el docente, como espacio
experimental para implementar y analizar el algoritmo A* junto con distintas funciones
heuristicas (nula, distancia Manhattan, distancia Euclidiana y heuristicas disenadas por el
grupo para los problemas de las esquinas y de la busqueda de alimentos). El proposito central no
es unicamente lograr que Pac-Man llegue al objetivo, sino comprender como la eleccion de la
heuristica afecta la cantidad de estados expandidos y el costo de la solucion encontrada,
preservando siempre la optimalidad mediante heuristicas admisibles y, cuando sea posible, consistentes.

El codigo base del proyecto (\texttt{pacman.py}, \texttt{search.py}, \texttt{searchAgents.py},
\texttt{util.py}, \texttt{layout.py}, entre otros) se mantiene sin modificar en su arquitectura;
el trabajo del grupo consiste en completar unicamente las funciones que la guia de laboratorio
senala como pendientes.
